\addcontentsline{toc}{chapter}{Abstract}\vspace{-1cm}
%Border
\begin{tikzpicture}[remember picture, overlay]
  \draw[line width = 4pt] ($(current page.north west) + (0.75in,-0.75in)$) rectangle ($(current page.south east) + (-0.75in,0.75in)$);
\end{tikzpicture}


Historical inscriptions and manuscripts across South Asia preserve invaluable linguistic, cultural, and archaeological information, yet many of these artefacts remain difficult to study because the available scans are degraded by noise, uneven illumination, fading ink, surface weathering, and background texture. These limitations reduce readability and make manual transcription slow, error-prone, and dependent on rare specialist expertise. This report addresses that gap by presenting a non-destructive digitisation pipeline that improves the visual quality of scanned artefacts and prepares them for searchable transcription without altering the original sources.

The work focuses on a staged workflow that begins with preprocessing and continues through AI-based enhancement, binarisation, OCR, and structured record generation. The pipeline combines image normalisation, denoising, contrast improvement, super-resolution, and adaptive thresholding to make degraded text more legible. OCR is then applied to the enhanced and binarised images to produce Unicode transcriptions, confidence values, and per-line text output, while the final record assembly step stores the results in a structured format suitable for archiving and later retrieval.

The implementation and evaluation use a software stack built around OpenCV, Pillow, Real-ESRGAN, Tesseract, EasyOCR, FastAPI, and a React-based web interface. Representative inscription and manuscript images from different material types, including stone, palm leaf, copper plate, paper, and cave or rock art, are used as test cases to verify the pipeline’s robustness across varied degradation patterns. The results show that the proposed workflow improves contrast, suppresses noise, and produces clearer character regions for transcription, making it a practical foundation for scalable digital preservation and future extension to translation and richer metadata extraction.

\pagebreak